\documentclass[11pt]{article}
\usepackage[margin=1in]{geometry}
\usepackage[T1]{fontenc}
\usepackage{lmodern,amsmath,amssymb,amsthm,mathtools,booktabs,longtable,microtype}
\usepackage[hidelinks]{hyperref}
\setlength{\emergencystretch}{3em}
\allowdisplaybreaks
\newtheorem{theorem}{Theorem}[section]
\newtheorem{lemma}[theorem]{Lemma}
\newtheorem{proposition}[theorem]{Proposition}
\title{Structured projector blocks and semantic bulk assembly\\
\large A conditional fixed-tape integer-multiplication bound}
\author{icekylinx\\\small With substantial OpenAI GPT-6 Astra and Codex assistance}
\date{8 October 2026}
\begin{document}
\maketitle
\begin{abstract}
A controlled rational basis with a fixed middle factor $I+J$, certified
contiguous pivot runs, and mixed Gaussian-dyadic centers give independent
bit and complex interfaces at dimensions $(32,30,36)$ and $(30,30,40)$.
We instantiate the attributed semantic precision, coordinate-routing,
and bulk resampling interfaces with these networks. Subject to the
retained analytic and fixed-tape framework, the result is
$T(n)=O(n(\log n)^{1-\kappa})$, with
$\kappa=2504683/200000000000=1.2523415\times10^{-5}$.
The finite certificates establish exact counts and strict inequalities;
the all-size conclusion also uses the written construction, analytic,
and tape proofs identified here.
\end{abstract}

\section{Scope, attribution, and inherited interfaces}
All computations use a fixed finite-alphabet multitape Turing machine.
Finite graphs, rational bases, address primes, tables, and tape counts
are fixed independently of multiplication input length. Constructive
setup, analytic recovery, and eventual thresholds remain part of the
retained conditional framework. This is not a formal verification or a
claim of an implemented complete multiplication machine.

The immediate repository predecessor is icekylinx's PR24 endpoint-gauge
construction \cite{parent} at \texttt{ed90fd9}, whose proof is \path{notes/endpoint-gauge-note.tex}.
Its partial-swap compiler, arbitrary-dirty-auxiliary scalar wrapper,
carrier semantics, controlled inner boundary prescription, A3 corner,
and arbitrary-integer-width bit rows are retained as specified below.
This follows icekylinx's PR10 batching and PR18 partial swaps, pinned at
\texttt{62691e395a0458ce089a1c7b5d89e74291e95e29} and
\texttt{f2ab41aebad47861caf6316282c1793e5513845e}, respectively.
The present finite refinements and their circuit-specific composition
are by icekylinx, with substantial OpenAI GPT-6 Astra/Codex assistance.

The new analytic composition uses Zhihao Chen's PR23 \cite{pr23}, pinned
at \texttt{661ebabc1076c9f525d7ce4967c876b203fa9827}, for the semantic
near-parent-child induction, literal scalar charge, simultaneous row
stock, and strict parameterization. The underlying semantic child idea,
arbitrary-coordinate routing, phase-cell inverse, and bulk locality and
tape constructions are due to the RaD project (hipotures) \cite{rad},
pinned at \texttt{45d9b60355872041f9275f77c057514def0d45bc} and exposed
through PR20. The translated complex sink is Zhihao Chen's PR21
\cite{pr21}, at \texttt{5ba6cf0bfb68f2be8d15610e7972207c50254d6a},
following eumemic's PR13 source frames and Chen's PR16 nested bases.
PR23's recorded OpenAI Codex assistance and RaD's recorded Codex and
GPT-6.1-Sol Ultra assistance retain their source attribution.

The foundation is Douglas Colkitt's compact-control framework, pinned
at \texttt{6e564879f51ae16f23d392e9e196c605f36d90df}, and the upstream
OpenAI manuscript at \texttt{adc7f1241b42e322a6451854ab7e4b4c146bf78a}.
Inherited contributions include Bortlesboat's alignment, eumemic's
complex producers and Gaussian resampling, dleen's retained totals and
stage sharing, Zhihao Chen's finite networks, Rohan Arun's related
geometry and matching, and Aurel Prosz's parameter analysis.
Their notices, licenses, recorded AI assistance, and the analytic debt
to David Harvey and Joris van der Hoeven are retained in \texttt{NOTICE}
and \texttt{SOURCES.json}. No new license condition or priority claim
is imposed. The source proof map appears in Section~\ref{sec:transfer}.

\section{Selected bit construction}
\input{notes/structured-bulk-bit.tex}
\section{Exact middle-axis pivot certificate}
\input{notes/structured-bulk-crt.tex}
\section{Complete rows and arbitrary bit widths}
The following inherited mixed-width construction uses the selected bit
child list above. Its triangular wrappers are the retained partial-swap
conjugation, and every contiguous pivot run is one child.
\input{notes/batched-bit-rows.tex}
\section{Mixed-center complex construction}
\input{notes/structured-bulk-complex.tex}
\section{Semantic precision, analytic transfer, and assembly}
\input{notes/structured-bulk-assembly.tex}

\begin{thebibliography}{9}
\bibitem{upstream} OpenAI. \emph{Integer multiplication below $n\log n$},
result family \#109. \url{https://github.com/openai/math}.
\bibitem{main} Douglas Colkitt et al. \emph{Integer multiplication bounds}.
\url{https://github.com/CrocSwap/integer-mult-bounds}.
\bibitem{parent} icekylinx. \emph{Endpoint gauges for integer multiplication}.
PR24, repository predecessor \texttt{ed90fd9}.
\url{https://github.com/CrocSwap/integer-mult-bounds/pull/24}.
\path{notes/endpoint-gauge-note.tex}.
\bibitem{pr10} icekylinx. Projector batching and controlled bases.
\url{https://github.com/CrocSwap/integer-mult-bounds/pull/10}.
\bibitem{pr18} icekylinx. Partial-swap frames and retained triples.
\url{https://github.com/CrocSwap/integer-mult-bounds/pull/18}.
\bibitem{pr21} Zhihao Chen (jacklightChen). Translated partial-swap frames.
\url{https://github.com/CrocSwap/integer-mult-bounds/pull/21}.
Retained proof: \path{references/semantic-bulk/pr21/translated-partial-complex.tex}.
\bibitem{pr23} Zhihao Chen (jacklightChen). Semantic precision and bulk
resampling. \url{https://github.com/CrocSwap/integer-mult-bounds/pull/23}.
Retained proof: \path{references/semantic-bulk/pr23/semantic-bulk-17-note.tex}.
\bibitem{rad} RaD project (hipotures). Analytic and tape constructions.
\url{https://github.com/hipotures/rad/tree/45d9b60355872041f9275f77c057514def0d45bc/research/integer-multiplication-bounds}.
Retained proofs: \path{references/semantic-bulk/rad20/reports/}.
\bibitem{bhp} R. C. Baker, G. Harman and J. Pintz.
The Difference Between Consecutive Primes, II.
\emph{Proc. London Math. Soc.} 83 (2001), 532--562, Theorem 1.
\end{thebibliography}
\end{document}
